Рассматривается система линейных алгебраических уравнений
\[
    Ax = b.
\]

Для применения итерационных методов систему необходимо привести к виду
\[
    x = Fx + c,
\]
где $F$ — матрица итерационного процесса, $c$ — вектор свободных членов.

\subsection*{Преобразование системы}

Пусть исходная система имеет вид
\[
    \sum_{j=1}^{n} a_{ij}x_j = b_i,
    \qquad i = 1,2,\ldots,n.
\]

Выразим из каждого уравнения неизвестную $x_i$:
\[
    x_i =
    -\sum_{\substack{j=1 \\ j \ne i}}^{n}
    \frac{a_{ij}}{a_{ii}}x_j
    +
    \frac{b_i}{a_{ii}}.
\]

Тогда элементы матрицы $F$ и вектора $c$ определяются следующим образом:
\[
f_{ij} =
\begin{cases}
0, & i = j, \\
-\dfrac{a_{ij}}{a_{ii}}, & i \ne j,
\end{cases}
\qquad
c_i = \frac{b_i}{a_{ii}}.
\]

После этого система принимает итерационный вид:
\[
    x = Fx + c.
\]

\subsection*{Сходимость метода простой итерации}

Метод простой итерации задается формулой
\[
    x^{(k)} = Fx^{(k-1)} + c,
    \qquad k = 1,2,\ldots
\]

В качестве начального приближения выбирается
\[
    x^{(0)} = c.
\]

Достаточным условием сходимости является выполнение неравенства
\[
    \|F\| < 1.
\]

Для оценки сходимости используется норма матрицы
\[
    \|F\| =
    \max_i \sum_{j=1}^{n} |f_{ij}|.
\]

На каждой итерации вычисляется абсолютная погрешность:
\[
    \Delta_k =
    \frac{\|F\|}{1 - \|F\|}
    \cdot
    \|x^{(k)} - x^{(k-1)}\|_{\infty}.
\]

Относительная погрешность определяется формулой
\[
    \delta_k =
    \frac{\Delta_k}{\|x^{(k)}\|_{\infty}}.
\]

Итерационный процесс продолжается до выполнения условия
\[
    \delta_k < 0.01.
\]

\subsection*{Диагональное преобладание}

Для обеспечения сходимости проверим условие диагонального преобладания:
\[
    |a_{ii}| >
    \sum_{\substack{j=1 \\ j \ne i}}^{n} |a_{ij}|.
\]

Для варианта $k = 22$:
\[
    \alpha = 2.2,
    \qquad
    \beta = 2.2.
\]

Тогда матрица коэффициентов имеет вид
\[
A =
\begin{pmatrix}
12.2 & -1.0 & 0.2 & 2.0 \\
1.0 & 9.8 & -2.0 & 0.1 \\
0.3 & -4.0 & 9.8 & 1.0 \\
0.2 & -0.3 & -0.5 & 5.8
\end{pmatrix}.
\]

Проверим строки матрицы:
\[
\begin{aligned}
12.2 &> 1.0 + 0.2 + 2.0 = 3.2, \\
9.8 &> 1.0 + 2.0 + 0.1 = 3.1, \\
9.8 &> 0.3 + 4.0 + 1.0 = 5.3, \\
5.8 &> 0.2 + 0.3 + 0.5 = 1.0.
\end{aligned}
\]

Во всех строках диагональный элемент по модулю больше суммы модулей остальных элементов строки. Следовательно, система уже имеет диагональное преобладание и дополнительного преобразования не требует.

\subsection*{Метод Зейделя}

Метод Зейделя отличается от метода простой итерации тем, что при вычислении очередной компоненты используются уже найденные значения текущей итерации.

Итерационная формула имеет вид
\[
x_i^{(k)}
=
\sum_{j=1}^{i-1} f_{ij}x_j^{(k)}
+
\sum_{j=i}^{n} f_{ij}x_j^{(k-1)}
+
c_i.
\]

Начальное приближение:
\[
    x^{(0)} = c.
\]

Критерий остановки:
\[
    \|x^{(k)} - x^{(k-1)}\|_{\infty} \le 10^{-4}.
\]

После выполнения данного условия итерационный процесс считается завершенным.